% GENERATED FILE. Edit canonical lessons, then run npm run generate:books.
% canonical-source-hash: fnv1a64:dff6019575e82f0e
% canonical-lessons: MR-C109-mhatara, MR-C109-athavda, MR-C109-mahina, MR-C109-varsa, MR-C109-tas

\chapter{Old, Week, Month, Year, Hour}
\label{ch:mr-old-week-month-year-hour}
\hlchaptermodality{109}

\begin{chapteropening}
\textbf{By the end of this chapter:} \emph{I can say old, a week, a month, a year and an hour.}

Complete the last lesson of chapter 109: I can say old, a week, a month, a year and an hour.
\end{chapteropening}

\section[\texorpdfstring{\mr{म्हातारा} (mhātārā)}{म्हातारा (mhātārā)}]{\mr{म्हातारा} (mhātārā) — old (people)}

\label{lesson:MR-C109-mhatara}



\begin{quote}
Before the new one: say the Marathi for pink, then the Marathi for young.
\end{quote}

\subsection*{You'll want to know: \mr{म्हातारा}}
\textbf{\mr{म्हातारा}} — \emph{mhātārā} — \textquotedblleft{}old (people)\textquotedblright{}.

\begin{cousinweb}
\textbf{\mr{म्हातारा}} (\emph{mhātārā}) is a tadbhava of Sanskrit \emph{mahattara}, \textquotedblleft{}greater\textquotedblright{}, which also meant \textquotedblleft{}chief, village eldest\textquotedblright{} (CDIAL 9932). CDIAL lists Marathi \emph{mahātārā}, \textquotedblleft{}old\textquotedblright{}.
\end{cousinweb}

\begin{grammarlens}[title={What you've built}]
One more word for this chapter.
\end{grammarlens}

\subsection*{Your turn}
\begin{itemize}
\raggedright
  \item \emph{Say it:} \emph{mhātārā}
  \item \emph{Say it:} \emph{mhātārā}, once more
  \item \emph{Say it:} say \emph{mhātārā}
  \item \emph{Recall:} say the Marathi for pink, then the Marathi for young, then say \emph{mhātārā} again
\end{itemize}

\begin{tcolorbox}[breakable,colback=teal!4,colframe=teal!35!black,title={Before you move on}]
What does \mr{म्हातारा} mean? (\textquotedblleft{}Old (people)\textquotedblright{}.) Say it once more.
\end{tcolorbox}
\section[\texorpdfstring{\mr{आठवडा} (āṭhavḍā)}{आठवडा (āṭhavḍā)}]{\mr{आठवडा} (āṭhavḍā) — a week}

\label{lesson:MR-C109-athavda}



\begin{quote}
Before the new one: say the Marathi for young, then the Marathi for old.
\end{quote}

\subsection*{You'll want to know: \mr{आठवडा}}
\textbf{\mr{आठवडा}} — \emph{āṭhavḍā} — \textquotedblleft{}a week\textquotedblright{}.

\begin{cousinweb}
\textbf{\mr{आठवडा}} (\emph{āṭhavḍā}) comes from \emph{*aṣṭavāra}, \textquotedblleft{}period of eight days\textquotedblright{}: \emph{aṣṭa}, \textquotedblleft{}eight\textquotedblright{}, and \emph{vāra}, \textquotedblleft{}turn, day\textquotedblright{} (CDIAL 940). CDIAL notes that the Marathi form was crossed with compounds of \emph{vāṭa}.
\end{cousinweb}

\begin{grammarlens}[title={What you've built}]
One more word for this chapter.
\end{grammarlens}

\subsection*{Your turn}
\begin{itemize}
\raggedright
  \item \emph{Say it:} \emph{āṭhavḍā}
  \item \emph{Say it:} \emph{āṭhavḍā}, once more
  \item \emph{Say it:} say \emph{āṭhavḍā}
  \item \emph{Recall:} say the Marathi for young, then the Marathi for old, then say \emph{āṭhavḍā} again
\end{itemize}

\begin{tcolorbox}[breakable,colback=teal!4,colframe=teal!35!black,title={Before you move on}]
What does \mr{आठवडा} mean? (\textquotedblleft{}A week\textquotedblright{}.) Say it once more.
\end{tcolorbox}
\section[\texorpdfstring{\mr{महिना} (mahinā)}{महिना (mahinā)}]{\mr{महिना} (mahinā) — a month}

\label{lesson:MR-C109-mahina}



\begin{quote}
Before the new one: say the Marathi for old, then the Marathi for a week.
\end{quote}

\subsection*{You'll want to know: \mr{महिना}}
\textbf{\mr{महिना}} — \emph{mahinā} — \textquotedblleft{}a month\textquotedblright{}.

\begin{cousinweb}
\textbf{\mr{महिना}} (\emph{mahinā}) is a loan, usually traced to Persian \emph{māhīna}, \textquotedblleft{}monthly\textquotedblright{}, built on \emph{māh}, \textquotedblleft{}moon, month\textquotedblright{}.
\end{cousinweb}

\begin{grammarlens}[title={What you've built}]
One more word for this chapter.
\end{grammarlens}

\subsection*{Your turn}
\begin{itemize}
\raggedright
  \item \emph{Say it:} \emph{mahinā}
  \item \emph{Say it:} \emph{mahinā}, once more
  \item \emph{Say it:} say \emph{mahinā}
  \item \emph{Recall:} say the Marathi for old, then the Marathi for a week, then say \emph{mahinā} again
\end{itemize}

\begin{tcolorbox}[breakable,colback=teal!4,colframe=teal!35!black,title={Before you move on}]
What does \mr{महिना} mean? (\textquotedblleft{}A month\textquotedblright{}.) Say it once more.
\end{tcolorbox}
\section[\texorpdfstring{\mr{वर्ष} (varṣa)}{वर्ष (varṣa)}]{\mr{वर्ष} (varṣa) — a year}

\label{lesson:MR-C109-varsa}



\begin{quote}
Before the new one: say the Marathi for a week, then the Marathi for a month.
\end{quote}

\subsection*{You'll want to know: \mr{वर्ष}}
\textbf{\mr{वर्ष}} — \emph{varṣa} — \textquotedblleft{}a year\textquotedblright{}.

\begin{cousinweb}
\textbf{\mr{वर्ष}} (\emph{varṣa}) is Sanskrit \emph{varṣa}, \textquotedblleft{}rain; year\textquotedblright{}, taken over whole as a learned loan, a tatsama (CDIAL 11392). It meant \textquotedblleft{}rain\textquotedblright{} in the Rigveda and \textquotedblleft{}year\textquotedblright{} later.
\end{cousinweb}

\begin{grammarlens}[title={What you've built}]
One more word for this chapter.
\end{grammarlens}

\subsection*{Your turn}
\begin{itemize}
\raggedright
  \item \emph{Say it:} \emph{varṣa}
  \item \emph{Say it:} \emph{varṣa}, once more
  \item \emph{Say it:} say \emph{varṣa}
  \item \emph{Recall:} say the Marathi for a week, then the Marathi for a month, then say \emph{varṣa} again
\end{itemize}

\begin{tcolorbox}[breakable,colback=teal!4,colframe=teal!35!black,title={Before you move on}]
What does \mr{वर्ष} mean? (\textquotedblleft{}A year\textquotedblright{}.) Say it once more.
\end{tcolorbox}
\section[\texorpdfstring{\mr{तास} (tās)}{तास (tās)}]{\mr{तास} (tās) — an hour}

\label{lesson:MR-C109-tas}



\begin{quote}
Before the new one: say the Marathi for a month, then the Marathi for a year.
\end{quote}

\subsection*{You'll want to know: \mr{तास}}
\textbf{\mr{तास}} — \emph{tās} — \textquotedblleft{}an hour\textquotedblright{}.

\begin{cousinweb}
\textbf{\mr{तास}} (\emph{tās}) is usually traced to Persian \emph{tās}, \textquotedblleft{}a metal bowl\textquotedblright{}. In Marathi it also names a gong, and the gong struck at each hour is taken to have given the hour its name. CDIAL has no entry for it.
\end{cousinweb}

\begin{grammarlens}[title={What you've built}]
One more word for this chapter.
\end{grammarlens}

\subsection*{Your turn}
\begin{itemize}
\raggedright
  \item \emph{Say it:} \emph{tās}
  \item \emph{Say it:} \emph{tās}, once more
  \item \emph{Say it:} say \emph{tās}
  \item \emph{Recall:} say the Marathi for a month, then the Marathi for a year, then say \emph{tās} again
\end{itemize}

\begin{tcolorbox}[breakable,colback=teal!4,colframe=teal!35!black,title={Before you move on}]
What does \mr{तास} mean? (\textquotedblleft{}An hour\textquotedblright{}.) Say it once more.
\end{tcolorbox}
